\section{Field by field}

The tables below spell out the headers. Offsets are in bytes from the start of the header.

\begin{table}[ht]
\centering
\begin{tabular}{rllp{7cm}}
\toprule
Offset & Size & Field & Meaning \\
\midrule
0 & 4 & \code{magic} & \code{0x43434242}, the bytes \code{42 42 43 43} on the wire \\
4 & 4 & \code{tag} & chosen by the host, echoed by the board \\
8 & 2 & \code{op} & operation code (Table~\ref{tab:ops}) \\
10 & 2 & \code{flags} & per operation flags, see below \\
12 & 4 & \code{id} & hash session id, or associated data length for \code{AEAD} \\
16 & 4 & \code{alg} & algorithm number (Table~\ref{tab:algs}) \\
20 & 4 & \code{keylen} & bytes of key that follow the header \\
24 & 4 & \code{ivlen} & bytes of IV that follow the key \\
28 & 4 & \code{datalen} & bytes of data that follow the IV \\
\bottomrule
\end{tabular}
\caption{Request header (32 bytes).}
\end{table}

\begin{table}[ht]
\centering
\begin{tabular}{rllp{7cm}}
\toprule
Offset & Size & Field & Meaning \\
\midrule
0 & 4 & \code{magic} & same constant \\
4 & 4 & \code{tag} & copy of the request's tag \\
8 & 4 & \code{status} & 0, or a negative \code{errno} as a 32 bit two's complement number \\
12 & 4 & \code{id} & new session id for \code{HASH\_NEW} and \code{HASH\_CLONE} \\
16 & 4 & \code{ivlen} & bytes of IV that follow (chained IV for ciphers) \\
20 & 4 & \code{datalen} & bytes of data that follow the IV \\
\bottomrule
\end{tabular}
\caption{Reply header (24 bytes).}
\end{table}

\begin{table}[ht]
\centering
\begin{tabular}{rl}
\toprule
Value & Operation \\
\midrule
0 & \code{PING} \\
1 & \code{HASH\_NEW} \\
2 & \code{HASH\_UPDATE} \\
3 & \code{HASH\_FINAL} \\
4 & \code{HASH\_CLONE} \\
5 & \code{HASH\_FREE} \\
6 & \code{CIPHER} \\
7 & \code{INFO} \\
8 & \code{AEAD} \\
\bottomrule
\end{tabular}
\caption{Operation codes.}\label{tab:ops}
\end{table}

\begin{table}[ht]
\centering
\begin{tabular}{rl}
\toprule
Value & Algorithm \\
\midrule
1 & SHA-1 \\
2 & SHA-256 \\
3 & HMAC-SHA-256 \\
4 & ECB-AES \\
5 & CBC-AES \\
6 & CTR-AES \\
7 & GCM-AES \\
\bottomrule
\end{tabular}
\caption{Algorithm numbers.}\label{tab:algs}
\end{table}

\paragraph{Flags.} Bit 0 means decrypt for \code{CIPHER} and \code{AEAD}. Bit 1 (\code{ONESHOT}) on \code{HASH\_NEW} asks for the digest of the attached data in the same command. For \code{AEAD}, bits 8 to 15 hold the tag length in bytes.

\section{Two messages in bytes}

\paragraph{INFO.} The first message the driver sends after probing a board is \code{INFO}, to check that the daemon is alive. With tag 42 it is a 32 byte header and nothing else:

\begin{lstlisting}[style=shell]
42 42 43 43   2a 00 00 00   07 00   00 00   00 00 00 00
magic         tag = 42      op=7    flags   id
00 00 00 00   00 00 00 00   00 00 00 00   00 00 00 00
alg           keylen        ivlen         datalen
\end{lstlisting}

The reply is a 24 byte header followed by the text that describes the board's algorithm drivers, for example \code{bbcryptod: sha1=... sha256=... cbc(aes)=cbc-aes-omap ...}. The driver stores the text and prints it when it announces the board.

\paragraph{One shot hash.} To hash the three bytes \code{abc} with SHA-256 in one command, the host sends operation 1 with flag \code{ONESHOT}, algorithm 2 and a data length of 3. The request is 35 bytes (a 32 byte header and the data). The board runs the hash and answers with a 24 byte header with status 0 and a data length of 32, followed by the 32 byte digest, in total 56 bytes. The digest of \code{abc} is the well known value

\begin{lstlisting}[style=shell]
ba7816bf 8f01cfea 414140de 5dae2223 b00361a3 96177a9c b410ff61 f20015ad
\end{lstlisting}

\paragraph{A request that needs padding.} A \code{PING} with 480 bytes of data is 32 + 480 = 512 bytes. A transfer of exactly 512 bytes would end in a full packet, so the sender adds one byte, and the board recognises and discards a following transfer of one byte.

\section{Messages for a hash session}

A hash of unknown length is a series of commands that share a session on the board (Figure~\ref{fig:hashseq}).

\begin{figure}[ht]
\centering
\begin{tikzpicture}[x=1cm,y=1cm]
  \draw[thick] (0,0) -- (0,-5.2); \node[above] at (0,0) {Host};
  \draw[thick] (7,0) -- (7,-5.2); \node[above] at (7,0) {Board};
  \draw[arr] (0,-0.5) -- node[lbl,above,sloped]{HASH\_NEW alg, key} (7,-0.8);
  \draw[arr] (7,-1.0) -- node[lbl,above,sloped]{reply id = 5} (0,-1.3);
  \draw[arr] (0,-1.7) -- node[lbl,above,sloped]{HASH\_UPDATE id 5, data} (7,-2.0);
  \draw[arr] (7,-2.2) -- node[lbl,above,sloped]{reply ok} (0,-2.5);
  \draw[arr] (0,-2.9) -- node[lbl,above,sloped]{HASH\_UPDATE id 5, data} (7,-3.2);
  \draw[arr] (7,-3.4) -- node[lbl,above,sloped]{reply ok} (0,-3.7);
  \draw[arr] (0,-4.1) -- node[lbl,above,sloped]{HASH\_FINAL id 5} (7,-4.4);
  \draw[arr] (7,-4.6) -- node[lbl,above,sloped]{reply digest, id freed} (0,-4.9);
\end{tikzpicture}
\caption{A streamed hash. Small inputs skip all of this with one shot hashing.}\label{fig:hashseq}
\end{figure}

The one shot flag exists because most hashes in practice are small, and a four message exchange for a 100 byte input would be dominated by round trips.

\section{Where the limits come from}

\begin{table}[ht]
\centering
\begin{tabularx}{\textwidth}{lrX}
\toprule
Limit & Value & Reason \\
\midrule
\code{BBC\_MAX\_DATA} & 32 KiB & one request must fit a bounded buffer on the board, and a bounded time on the wire \\
\code{BBC\_MAX\_KEY} & 128 bytes & AES keys are at most 32 bytes. HMAC keys can be as long as the hash block (64 bytes), and the extra headroom is cheap \\
\code{BBC\_MAX\_IV} & 16 bytes & the AES block size \\
Window & 1 to 8, default 4 & more in flight than the board's receive slots (12) and reply cache (16) can cover would break the retransmission guarantee \\
Retries & 4 sends, 1 s each & long enough for a busy daemon, short enough that a dead one is noticed \\
Reply cache & 16 entries & must exceed the window, and must exceed the writer queue (8) so a cached reply is not overwritten while it is still being sent \\
Hash sessions & 256 & each is an open socket on the board; the oldest is evicted if the table is full \\
\bottomrule
\end{tabularx}
\caption{Limits and why they have those values.}
\end{table}

\section{Status codes}

The status field in a reply is an ordinary Linux error number, negated. The driver passes it straight on to the caller, which is why a forged GCM tag reaches a user program as \code{EBADMSG} through the socket. The codes that occur are \code{-EINVAL} (bad lengths or parameters), \code{-ENOENT} (unknown hash session), \code{-EOPNOTSUPP} (unknown algorithm), \code{-ENOSYS} (unknown operation), \code{-ENOMEM}, and \code{-EBADMSG} for a failed tag check. Errors from the transport itself (\code{-ENODEV}, \code{-ETIMEDOUT}, \code{-EPROTO}, \code{-EMSGSIZE}) are generated on the host and never appear in a reply.
